% !TeX root = ../main.tex
\chapter*{ĐẶT VẤN ĐỀ}
\addcontentsline{toc}{chapter}{ĐẶT VẤN ĐỀ}
\markboth{ĐẶT VẤN ĐỀ}{ĐẶT VẤN ĐỀ}

\section*{1. Lý do chọn đề tài}

Trong bối cảnh chuyển đổi số, một tổ chức thường vận hành đồng thời nhiều ứng
dụng và dịch vụ khác nhau: thư điện tử, hệ thống quản trị nội bộ, cổng thông
tin, các ứng dụng nghiệp vụ, dịch vụ đám mây, v.v. Nếu mỗi ứng dụng tự quản lý
tài khoản riêng, người dùng phải ghi nhớ nhiều bộ tài khoản/mật khẩu, dẫn tới
việc tái sử dụng mật khẩu, đặt mật khẩu yếu và gia tăng nguy cơ bị tấn công.
Về phía quản trị, việc cấp phát, thu hồi và đồng bộ tài khoản thủ công trên
nhiều hệ thống vừa tốn thời gian, vừa khó kiểm soát và dễ sai sót.

Một hệ thống xác thực tập trung (xây dựng quanh một
\textit{Identity Provider} -- IdP) cho phép nhiều ứng dụng cùng tin cậy vào
một nguồn định danh duy nhất, hỗ trợ đăng nhập một lần (Single
Sign-On -- SSO), tập trung hóa chính sách bảo mật như xác thực đa yếu tố
(MFA) và cung cấp một điểm kiểm soát, giám sát duy nhất. Đây là mô hình đã
được các tổ chức lớn áp dụng rộng rãi.

Song song đó, bản thân dịch vụ xác thực là một dịch vụ có tính trọng yếu (trust-critical): nó có trạng thái (cơ sở dữ liệu, phiên), yêu cầu bảo
mật cao và là điểm truy cập bắt buộc của mọi ứng dụng. Việc triển khai và vận
hành dịch vụ này một cách thủ công dễ dẫn tới sai sót cấu hình, khó
tái lập môi trường, thời gian khắc phục sự cố dài và khó mở rộng. Do đó, tự
động hóa quá trình triển khai và vận hành dịch vụ xác thực là một bài toán có
ý nghĩa thực tiễn rõ ràng.

Từ hai nhận định trên, chuyên đề lựa chọn hướng tiếp cận kết hợp: xây dựng một
hệ thống xác thực tập trung dựa trên Keycloak, đồng thời nghiên cứu và xây
dựng quy trình triển khai tự động cho chính dịch vụ đó.

\section*{2. Mục tiêu của chuyên đề}

Chuyên đề hướng tới các mục tiêu sau:

\begin{itemize}
  \item \textbf{Mục tiêu 1:} Nghiên cứu cơ sở lý thuyết về xác thực tập trung
  và các giao thức liên quan (OAuth 2.0, OpenID Connect), từ đó thiết kế một
  hệ thống định danh tập trung hỗ trợ SSO dựa trên Keycloak.
  \item \textbf{Mục tiêu 2:} Xây dựng quy trình quản lý cấu hình định danh như mã nguồn (Configuration-as-Code) và tự động hóa việc triển khai
  dịch vụ bằng các công cụ CI/CD, IaC và GitOps.
  \item \textbf{Mục tiêu 3:} Đánh giá hệ thống ở các khía cạnh chức năng, khả
  năng tự động hóa và bảo mật; đề xuất hướng phát triển thành đồ án tốt nghiệp.
\end{itemize}

\section*{3. Đối tượng và phạm vi nghiên cứu}

\textbf{Đối tượng nghiên cứu:} các giao thức xác thực/ủy quyền (OAuth 2.0,
OIDC), nền tảng Keycloak, công nghệ container hóa và các phương pháp triển
khai tự động.

\textbf{Phạm vi nghiên cứu:} chuyên đề tập trung vào:
\begin{itemize}
  \item Xây dựng IdP tập trung với Keycloak, PostgreSQL và một ứng dụng demo
  đóng vai trò Relying Party sử dụng luồng \textit{Authorization Code} kèm
  PKCE.
  \item Quản lý cấu hình định danh (realm, client, role, chính sách MFA) dưới
  dạng mã nguồn.
  \item Thiết kế và bước đầu hiện thực quy trình triển khai tự động trên nền
  tảng Kubernetes (k3s) theo mô hình GitOps.
\end{itemize}

Trong khuôn khổ thời gian của chuyên đề, một số hạng mục nâng cao (đa vùng,
cân bằng tải nâng cao, kiểm thử hiệu năng quy mô lớn) được trình bày ở mức
thiết kế và định hướng, làm cơ sở mở rộng thành đồ án tốt nghiệp.

\section*{4. Phương pháp nghiên cứu}

\begin{itemize}
  \item \textbf{Phương pháp nghiên cứu tài liệu:} đọc các đặc tả chuẩn
  (RFC 6749, RFC 7636, OIDC Core), tài liệu Keycloak và tài liệu Kubernetes.
  \item \textbf{Phương pháp phân tích -- thiết kế:} phân tích yêu cầu, mô hình
  hóa kiến trúc (mô hình C4), vẽ biểu đồ tuần tự cho luồng xác thực.
  \item \textbf{Phương pháp thực nghiệm:} triển khai hệ thống thực tế, cấu
  hình, kiểm thử chức năng và đo lường.
  \item \textbf{Phương pháp đánh giá:} đánh giá theo các tiêu chí chức năng,
  hiệu năng và bảo mật (threat model STRIDE).
\end{itemize}

\section*{5. Ý nghĩa của chuyên đề}

Về mặt khoa học, chuyên đề hệ thống hóa kiến thức về định danh tập trung và
vận hành dịch vụ trọng yếu. Về mặt thực tiễn, kết quả của chuyên đề có thể
được áp dụng để xây dựng hệ thống SSO nội bộ cho tổ chức, giảm chi phí quản
trị và tăng cường bảo mật, đồng thời là nền tảng để phát triển thành đồ án tốt
nghiệp ở mức độ hoàn chỉnh hơn.

\section*{6. Đóng góp của chuyên đề}

Chuyên đề có một số đóng góp chính: (1) hệ thống hóa cơ sở lý thuyết và so sánh
các giải pháp định danh phổ biến; (2) thiết kế kiến trúc ba tầng kết hợp dịch vụ
định danh với quy trình tự động hóa; (3) triển khai thực nghiệm Keycloak cùng
cấu hình định danh như mã nguồn và ứng dụng demo SSO; và (4) xây dựng khung
quy trình CI/CD, GitOps, giám sát cùng phân tích bảo mật theo STRIDE.

\section*{7. Kế hoạch thực hiện}

Bảng~\ref{tab:kehoach} trình bày kế hoạch thực hiện chuyên đề theo giai đoạn.

\begin{table}[H]
\centering
\caption{Kế hoạch thực hiện chuyên đề theo giai đoạn}
\label{tab:kehoach}
\small
\begin{tabular}{c p{8cm} c}
\toprule
\textbf{Giai đoạn} & \textbf{Nội dung} & \textbf{Thời gian}\\
\midrule
1 & Nghiên cứu lý thuyết và khảo sát công nghệ & Tuần 1--2\\
2 & Phân tích yêu cầu và thiết kế kiến trúc & Tuần 3--4\\
3 & Triển khai dịch vụ định danh và ứng dụng demo & Tuần 5--7\\
4 & Tự động hóa và quản lý cấu hình như mã nguồn & Tuần 8--10\\
5 & Đánh giá và hoàn thiện báo cáo & Tuần 11--14\\
\bottomrule
\end{tabular}
\end{table}

\section*{8. Bố cục báo cáo}

Ngoài phần Đặt vấn đề và Kết luận, báo cáo được tổ chức thành bốn chương:
\begin{itemize}
  \item \textbf{Chương 1} trình bày cơ sở lý thuyết về xác thực tập trung,
  giao thức OAuth 2.0/OIDC, Keycloak và triển khai tự động.
  \item \textbf{Chương 2} phân tích yêu cầu và thiết kế kiến trúc hệ thống.
  \item \textbf{Chương 3} trình bày quá trình triển khai và cấu hình hệ thống.
  \item \textbf{Chương 4} đánh giá kết quả và thảo luận.
\end{itemize}
